% 13-borrows.tex — Chapter 13: Borrow Checking and Regions
\chapter{Borrow Checking and Regions}
\label{chap:13-borrows}
\index{borrow checking}
\index{regions}
\index{two-phase borrows}

\pnum
The borrow checker ensures that references do not outlive their owners and that
the aliasing invariant (one mutable borrow \emph{or} many shared borrows) is
maintained.  It is implemented in \srcref{src/sema/region.h} and is invoked from
the linearity checker (\srcref{src/sema/linearity.h}).

% -------------------------------------------------------------------------
\section{Regions}
\label{sec:borrow-regions}
\index{Region struct}
% -------------------------------------------------------------------------

\pnum
A \cname{Region} represents a lexical scope boundary.  Regions are organized
in a parent chain from inner to outer:

\begin{structdoc}{Region}{src/sema/region.h:38}
\cname{id}     & \cname{int}      & Unique identifier (monotonically increasing). \\
\cname{depth}  & \cname{int}      & Nesting depth (0 = function scope). \\
\cname{parent} & \cname{Region *} & Enclosing region (NULL for the outermost scope). \\
\end{structdoc}

\pnum
\cname{region\_contains(outer, inner)} walks the parent chain of \texttt{inner}
to check whether \texttt{outer} appears in it.  If so, \texttt{inner} is a
sub-scope of \texttt{outer} and will terminate before \texttt{outer}.

\pnum
The invariant checked by the borrow checker is: the borrow region must be
contained within the owner region.  That is, \cname{region\_contains(owner\_region, borrow\_region)} must hold, guaranteeing the owner outlives the borrow.

% -------------------------------------------------------------------------
\section{Borrow entries}
\label{sec:borrow-entries}
\index{BorrowEntry}
% -------------------------------------------------------------------------

\pnum
An active borrow is recorded as a \cname{BorrowEntry}:

\begin{structdoc}{BorrowEntry}{src/sema/region.h:78}
\cname{var}           & \cname{Id *}        & The borrowed variable. \\
\cname{owner\_var}    & \cname{Id *}        & Direct owner of the borrow. \\
\cname{root\_owner}   & \cname{Id *}        & Ultimate non-reference owner (for transitive chains). \\
\cname{mode}          & \cname{OwnershipMode} & \cname{MODE\_SHARED} or \cname{MODE\_MUTABLE}. \\
\cname{phase}         & \cname{BorrowPhase}  & \cname{BORROW\_RESERVED} or \cname{BORROW\_ACTIVE}. \\
\cname{borrow\_region}& \cname{Region *}    & Scope where the borrow is used. \\
\cname{owner\_region} & \cname{Region *}    & Scope where the owner is defined. \\
\cname{is\_temporary} & \cname{bool}        & Expires at end of current statement. \\
\cname{last\_use\_stmt\_idx}& \cname{int}   & NLL: last statement index where binding is used. \\
\end{structdoc}

% -------------------------------------------------------------------------
\section{Two-phase borrows}
\label{sec:borrow-twophase}
\index{two-phase borrows}
\index{BORROW\_RESERVED}
\index{BORROW\_ACTIVE}
% -------------------------------------------------------------------------

\pnum
Lain implements two-phase borrows to permit calling a method on a container while
simultaneously passing container elements as arguments.  A borrow begins in the
\cname{BORROW\_RESERVED} phase, which allows shared reads but does not yet
activate the exclusivity constraint.  After all call arguments have been
evaluated, the borrow is promoted to \cname{BORROW\_ACTIVE} via
\cname{borrow\_promote\_to\_active()}.

\pnum
The promotion sequence in the expression handler for \cname{EXPR\_CALL}:
\begin{enumerate}
\item Register a reserved borrow for the implicit receiver argument.
\item Evaluate all call arguments (with the borrow in \cname{BORROW\_RESERVED}).
\item Call \cname{borrow\_promote\_to\_active()} to activate the borrow.
\end{enumerate}

\pnum
This follows the model described in \textit{Two-Phase Borrows} (Matsakis 2018).
The key insight is that a reserved borrow does not conflict with shared reads
of the same variable, so the argument evaluation (which may read the container)
is allowed before the borrow becomes exclusive.

% -------------------------------------------------------------------------
\section{Aliasing rules}
\label{sec:borrow-aliasing}
\index{aliasing rules}
% -------------------------------------------------------------------------

\pnum
The \cname{BorrowTable} enforces the following aliasing rules at every borrow
registration:

\begin{enumerate}
\item \textbf{No mutable borrow while any borrow is active.}  If the requested
  mode is \cname{MODE\_MUTABLE} and any active borrow of the same root owner
  exists, the operation is rejected.
\item \textbf{No shared borrow while a mutable borrow is active.}  If the
  requested mode is \cname{MODE\_SHARED} and an active mutable borrow of the
  same root owner exists, the operation is rejected.
\end{enumerate}

\pnum
Temporary borrows (\cname{is\_temporary = true}) are released at the end of
each statement.  Non-temporary borrows are released when their binding goes
out of scope.

% -------------------------------------------------------------------------
\section{NLL last-use optimization}
\label{sec:borrow-nll}
\index{NLL}
% -------------------------------------------------------------------------

\pnum
The \cname{last\_use\_stmt\_idx} field supports Non-Lexical Lifetime (NLL)
shortening.  The use-analysis pre-pass (\srcref{src/sema/use\_analysis.h})
assigns to each borrow the index of the last top-level statement that uses the
binding.  The borrow checker can then expire a borrow as soon as it is no longer
used, even before the lexical scope boundary is reached.  This allows code like:

\begin{laincode}
var x = acquire_resource()
use x as r
// r's borrow expires here (last use above)
// x can now be moved below:
f(mov x)
\end{laincode}
